\subsection{Degree of an extension}

Let A be an algebra over a field K. It is in particular a vector space over K; the dimension of this vector space is called the degree of A over K and is written \([A:K]\) (II, p.~293). By definition \([A:K]\) is thus the cardinal of any basis of A over K. This definition applies in particular to the case of extensions of K.

An extension of degree 1 is trivial. An extension of degree 2 (resp. 3 etc.) is called quadratic (resp. cubic etc.). An extension of finite degree is sometimes called, by abuse of language, a finite extension.

THEOREM 1. --- Let E be an extension of K and A an algebra over E. Then we have \([A:K]=[A:E]\) . \([E:K]\) . In particular, if F is an extension of E, we have

\[
[ \mathbf {F}: \mathbf {K} ] = [ \mathbf {F}: \mathbf {E} ]. [ \mathbf {E}: \mathbf {K} ].\tag{1}
\]

The theorem is just a special case of II, p.~222, Prop. 25; more precisely, if \((a_{\lambda})_{\lambda \in L}\) is a basis of A over E and \((b_{\mu})_{\mu \in M}\) a basis of E over K, then the family \((a_{\lambda}b_{\mu})_{(\lambda ,\mu)\in L\times M}\) is a basis of A over K.

COROLLARY 1. --- Let K, E, F be three fields such that \(K \subset E \subset F\) and \([F:K]\) is finite. Then the degrees \([E:K]\) and \([F:E]\) are divisors of \([F:K]\) .

If the degree \([F:K]\) is prime, there is thus no subextension of F other than K and F. But note that when \([F:K]\) is not prime, there is not necessarily a subextension of F other than K and F (cf.~V, p.~146, Exercise 1).

COROLLARY 2. --- Let K, E and F be three fields such that \(K \subset E \subset F\) . Suppose that \([F:K]\) is finite, then the relation \([E:K] = [F:K]\) is equivalent to E = F and the relation \([F:E] = [F:K]\) is equivalent to E = K.

For if L is an extension field of \(L'\) then \([L:L'] = 1\) is equivalent to \(L' = L\) .

PROPOSITION 1. --- Let A be an algebra of finite degree over a field K. If an element \(a \in \mathbf{A}\) is not a left (resp. right) divisor of zero in A, then it is invertible in A.

For the vector space A over K is of finite dimension, by hypothesis, and the linear mapping \(x \mapsto ax\) (resp. \(x \mapsto xa\) ) of A into A is injective; it is therefore bijective (II, p.~298, Cor.) and hence (I, p.~16, Remark) a is invertible in A.

COROLLARY. --- Let A be a commutative algebra of finite degree over a field K. If A is an integral domain, then it is a field.

